% figure 
\begin{figure}[!htb] 
	\tikzsetnextfilename{fig_KBM_typical_impulse_response}
	\centering
		\resizebox{\textwidth}{!}{
			\begin{tikzpicture} [
				/pgf/declare function={
					omega_n1 = 13.8744;
					xi1 = 0.5663;
					K1 = 539;
					omega_n2 = 21.6025;
					xi2 = 0.4629;
					K2 = 280;
					}
				]
			\centering
			%
			\begin{axis}[%
				name=ax1,
				width=.95\linewidth,
				height=.3\linewidth,
				scale only axis,   
				xlabel={temps (\si{\second})},
				ylabel={Position (\si{\centi\meter})},
				yticklabel={\pgfmathparse{\tick*100}\pgfmathprintnumber{\pgfmathresult}},
				axis x line = bottom,
				axis y line = left, 
				ylabel near ticks,
				]
				%
				\addplot [
					color=Orient,
					smooth,
					domain=0:1,
					samples=1000,
					]
				{omega_n1/(sqrt(1-xi1^2))*exp(-xi1*omega_n1*x)*sin(deg(omega_n1*sqrt(1-xi1^2)*x))/K1};
				%
				\addplot [
					color=Shiraz,
					smooth,
					domain=0:1,
					samples=1000,
				]
				{omega_n1/(sqrt(1-xi2^2))*exp(-xi2*omega_n2*x)*sin(deg(omega_n2*sqrt(1-xi2^2)*x))/K2};
				\addlegendentry{$(539,44,2.8)$};	
				\addlegendentry{$(280,12,0.6)$};	
			\end{axis}
			
			\end{tikzpicture}
		}	
	\caption{Réponse impulsionnelle en position du modèle viscoélastique à une impulsion d'amplitude unitaire (\SI{1}{\newton}), pour les valeurs extrêmes des paramètres identifiés par \textcite{erden_hand_2015}.}
	\label{fig_KBM_typical_impulse_response}
\end{figure}
	%%
	%